% 页面与正文
\usepackage[
  top=2.3cm,
  bottom=2.3cm,
  left=2.45cm,
  right=2.45cm,
  headheight=22pt
]{geometry}
\usepackage{setspace}
\setstretch{1.24}
\setlength{\parindent}{2em}
\setlength{\parskip}{0.25em}

% 颜色
\usepackage{xcolor}
\definecolor{SoftBlue}{RGB}{76,112,150}
\definecolor{SoftGray}{RGB}{245,247,250}
\definecolor{BlockBlue}{RGB}{220,232,245}
\definecolor{BlockBlueDark}{RGB}{200,220,240}
\definecolor{BaseGray}{RGB}{135,145,158}
\definecolor{COPFBlue}{RGB}{60,102,170}
\definecolor{WarmOrange}{RGB}{210,126,71}
\definecolor{SoftGreen}{RGB}{83,145,126}
\definecolor{DarkText}{RGB}{28,31,35}
\definecolor{LightTan}{RGB}{247,243,236}
\color{DarkText}

% 数学、表格与图片
\usepackage{amsmath,amssymb,mathtools}
\usepackage{graphicx}
\graphicspath{{figures/}}
\usepackage{booktabs}
\usepackage{tabularx}
\usepackage{longtable}
\usepackage{array}
\usepackage{multirow}
\newcolumntype{Y}{>{\raggedright\arraybackslash}X}
\newcolumntype{C}{>{\centering\arraybackslash}X}

% TikZ 与结构示意图
\usepackage{tikz}
\usetikzlibrary{
  arrows.meta,
  positioning,
  calc,
  fit,
  backgrounds,
  shapes.geometric
}
\usepackage{pgfplots}
\pgfplotsset{compat=1.18}
\tikzset{
  quantArrow/.style={
    -{Latex[length=2.4mm]},
    line width=1.15pt,
    draw=SoftBlue
  },
  quantSoftArrow/.style={
    -{Latex[length=2.1mm]},
    line width=0.9pt,
    draw=gray!55
  },
  quantCard/.style={
    rounded corners=5pt,
    draw=SoftBlue!38,
    fill=SoftGray,
    line width=0.8pt,
    inner sep=6pt
  },
  quantBlueCard/.style={
    rounded corners=5pt,
    draw=COPFBlue!55,
    fill=BlockBlue,
    line width=0.95pt,
    inner sep=6pt
  },
  quantOrangeCard/.style={
    rounded corners=5pt,
    draw=WarmOrange!70!black,
    fill=WarmOrange!8,
    line width=0.95pt,
    inner sep=6pt
  },
  quantGreenCard/.style={
    rounded corners=5pt,
    draw=SoftGreen!70!black,
    fill=SoftGreen!8,
    line width=0.95pt,
    inner sep=6pt
  }
}

% 语义提示框
\usepackage[most]{tcolorbox}
\tcbset{
  lecturebox/.style={
    enhanced,
    breakable,
    boxrule=0.8pt,
    arc=3pt,
    left=7pt,
    right=7pt,
    top=6pt,
    bottom=6pt,
    before skip=10pt,
    after skip=10pt,
    fonttitle=\bfseries
  }
}
\newtcolorbox{principlebox}[1][]{
  lecturebox,
  colback=BlockBlue!68,
  colframe=SoftBlue,
  coltitle=DarkText,
  title={通用原则},
  #1
}
\newtcolorbox{teachingbox}[1][]{
  lecturebox,
  colback=SoftGray,
  colframe=BaseGray!55,
  coltitle=DarkText,
  title={教学简化},
  #1
}
\newtcolorbox{warningbox}[1][]{
  lecturebox,
  colback=WarmOrange!8,
  colframe=WarmOrange,
  coltitle=DarkText,
  title={常见错误},
  #1
}
\newtcolorbox{confirmbox}[1][]{
  lecturebox,
  colback=WarmOrange!6,
  colframe=WarmOrange!85!black,
  coltitle=DarkText,
  title={需实际确认},
  #1
}
\newtcolorbox{practicebox}[1][]{
  lecturebox,
  colback=SoftGreen!8,
  colframe=SoftGreen,
  coltitle=DarkText,
  title={正确实践},
  #1
}
\newtcolorbox{casebox}[1][]{
  lecturebox,
  colback=LightTan,
  colframe=WarmOrange!55,
  coltitle=DarkText,
  title={贯穿案例},
  #1
}
\newtcolorbox{outputbox}[1][]{
  lecturebox,
  colback=SoftGreen!7,
  colframe=SoftGreen!85!black,
  coltitle=DarkText,
  title={本章产出},
  #1
}
\newtcolorbox{quizbox}[1][]{
  lecturebox,
  colback=SoftGray,
  colframe=SoftBlue!65,
  coltitle=DarkText,
  title={章节自测},
  #1
}

% Python 代码
\usepackage{listings}
\lstdefinestyle{pythonstyle}{
  language=Python,
  basicstyle=\ttfamily\small,
  backgroundcolor=\color{SoftGray},
  keywordstyle=\color{COPFBlue}\bfseries,
  commentstyle=\color{SoftGreen!70!black},
  stringstyle=\color{WarmOrange!85!black},
  numbers=left,
  numberstyle=\tiny\color{BaseGray},
  numbersep=8pt,
  frame=single,
  rulecolor=\color{BaseGray!35},
  breaklines=true,
  showstringspaces=false,
  tabsize=4,
  columns=fullflexible,
  captionpos=b
}
\lstset{style=pythonstyle}

% 列表与图注
\usepackage{enumitem}
\setlist[itemize]{
  leftmargin=2em,
  itemsep=0.25em,
  topsep=0.35em
}
\setlist[enumerate]{
  leftmargin=2.2em,
  itemsep=0.3em,
  topsep=0.35em
}
\usepackage{caption}
\captionsetup{
  font=small,
  labelfont=bf,
  labelsep=quad
}

% 章节标题
\ctexset{
  chapter={
    format=\raggedright\Huge\bfseries\color{DarkText},
    nameformat=\color{SoftBlue},
    numberformat=\color{SoftBlue},
    titleformat=\color{DarkText},
    beforeskip=-5pt,
    afterskip=24pt
  },
  section={
    format=\Large\bfseries\color{SoftBlue},
    beforeskip=1.4ex,
    afterskip=0.8ex
  },
  subsection={
    format=\large\bfseries\color{DarkText},
    beforeskip=1.2ex,
    afterskip=0.55ex
  }
}
\newcommand{\keyterm}[1]{\textbf{\color{SoftBlue}#1}}
\newcommand{\riskterm}[1]{\textbf{\color{WarmOrange!90!black}#1}}
\newcommand{\goodterm}[1]{\textbf{\color{SoftGreen!80!black}#1}}

% 页眉页脚
\usepackage{fancyhdr}
\pagestyle{fancy}
\fancyhf{}
\fancyhead[L]{\small\color{BaseGray}\leftmark}
\fancyhead[R]{\small\color{BaseGray}量化研究入门}
\fancyfoot[C]{\small\color{BaseGray}\thepage}
\renewcommand{\headrulewidth}{0.4pt}
\renewcommand{\headrule}{
  \hbox to\headwidth{
    \color{SoftBlue!45}\leaders\hrule height \headrulewidth\hfill
  }
}

% 链接与交叉引用
\usepackage[
  colorlinks=true,
  linkcolor=SoftBlue,
  urlcolor=SoftBlue,
  citecolor=SoftGreen!70!black
]{hyperref}
\usepackage[nameinlink,noabbrev]{cleveref}
